\section{The classification and its scope}\label{sec:introduction}

An integral periodic orbit of one polynomial automorphism is a different
object from a finite orbit of the group containing that automorphism.
For Fricke cubics this distinction matters: the geometry of the Vieta
group provides reduction methods, but the periodic points of one
prescribed word must still be exhausted in that word's own time.
We study the word containing the three Vieta involutions in a fixed order.
The result separates its unbounded integral periodic families from a
bounded, effectively decidable remainder, with one bound for all levels.

Fix an arbitrary ordered triple \((A,B,C)\in\Z^3\), and define
\begin{equation}\label{eq:invariant}
 K(x,y,z)=x^2+y^2+z^2-xyz-Ax-By-Cz.
\end{equation}
The Vieta involutions are
\begin{align}
 s_x(x,y,z)&=(A+yz-x,y,z),\notag\\
 s_y(x,y,z)&=(x,B+xz-y,z),\label{eq:involutions}\\
 s_z(x,y,z)&=(x,y,C+xy-z).\notag
\end{align}
Our native map is
\begin{equation}\label{eq:native}
 T=s_z\circ s_y\circ s_x,
\end{equation}
so the rightmost factor acts first. One application of \(T\), not one
individual involution, is one unit of time. Each factor preserves
\(K\), and each is a polynomial involution of \(\Z^3\). In particular,
no point is removed when a level \(K=D\) is singular.

Write \(\Per_\Z(T)=\{P\in\Z^3:T^nP=P\text{ for some }n\ge1\}\), and set
\begin{equation}\label{eq:bounds}
 \begin{gathered}
 H=\max(|A|,|B|,|C|),\qquad B_0=H+1,\qquad R=100B_0,\\
 Q_R=[-R,R]^3\cap\Z^3,\qquad N_R=(2R+1)^3.
 \end{gathered}
For a periodic point, its \emph{least period} is its first positive
return time. A \emph{certified return time} need not be the least period.

\begin{theorem}[All-coefficient, all-level integral classification]
\label{thm:main}
For every ordered \((A,B,C)\in\Z^3\), there is an explicit construction
of a union \(\lineset\) of at most 27 affine integer lines such that:
\begin{enumerate}
 \item Each line has a parametrization \(P(t)=p+tv\), with
 \(p,v\in\Z^3\), containing a coordinate of slope \(1\) or \(-1\).
 It consists entirely of periodic points and has a certified native
 return time \(r\le54\): \(T^rP(t)=P(t)\) for all \(t\in\Z\).
 \item If \(\core\) denotes the cyclic vertices of the exact partial
 map \(T:Q_R\dashrightarrow Q_R\), then
 \begin{equation}\label{eq:exhaustive_union}
    \Per_\Z(T)=\lineset\cup\core.
 \end{equation}
 Every periodic orbit not contained in \(\lineset\) lies entirely in
 \(Q_R\).
 \item A finite parameterwise procedure determines coincident lines,
 all integral intersections, the generic least period and every
 exceptional parameter on each line, and all remaining finite cycles.
 Thus all least periods and their coexistence are effective. Every
 native least period is at most \(\max(54,N_R)\).
 \item For every \(D\in\Z\), the level \(K=D\) contains at most
 \(54+N_R\) periodic integer points. Their exact least periods and
 their ordinary primitive-cycle counts are effectively determined.
\end{enumerate}
The construction uses the 27-state graph of Section~\ref{sec:graph}
and its affine-return test. Neither the numerical bounds nor the
resulting algorithmic complexity are asserted optimal.
\end{theorem}

The line part is a finite list of free integer parametrizations. The
core part is a proved terminating rule for the given coefficients, not
a table that has been computed for every coefficient triple. There is
no empirical parameter window or period cutoff in
Theorem~\ref{thm:main}. Its proof is given completely below.

\subsection{Related results and the contribution retained here}
\label{subsec:ownership}

Shin's reduction theory uses height graphs and low-coordinate Vieta
identities to study integral group orbits
\citep[Theorems 1.1--1.3 and \S2]{shin2026character}.
In particular, the orbit-equivalence algorithms concern the full
Vieta group and the mapping class group. Remark~3.1 in the accepted
manuscript explicitly removes the boundary-trace relations between
the coefficients. Independent coefficients and the height-graph and
low-coordinate mechanisms are therefore prior work, not an extension
claimed here. The integral change \((x,y,z)\mapsto(-x,y,z)\) identifies
our sign convention with Shin's, with parameters
\((\alpha_1,\alpha_2,\alpha_3,\beta)=(-A,B,C,D)\), and conjugates
the respective ordered involutions. The additional assertion proved
here is the fixed-word, level-uniform line/core exhaustion, not
decidability of whole-group orbit equivalence.

Cantat studies positive-entropy automorphisms on the same cubic family
\citep[Proposition 2.2 and \S3]{cantat2009bers}.
The fixed-fibre escape neighborhoods in that analysis imply compact
containment of periodic points, and hence integral finiteness on a fixed
fibre. This consequence is already available and is not our novelty
claim. It does not furnish the particular level-independent radius or
integer line construction used here. Our periodic lines cross invariant
levels; they are not invariant curves inside a fixed cubic surface, so
they do not conflict with Cantat's Corollary~3.4.

The equal-forcing slice is also prior work. The local working manuscript
C421 classifies all integer periodic orbits of
\(F_a(x,y,z)=(y,z,yz+a-x)\), including its exceptional families and
computer-assisted finite core \citep{c4212026integral}.
When \(A=B=C=a\), our map is exactly \(F_a^3\). That slice is not counted
as a new classification here, and its finite certificate is not an
input to our proof. The phase-dependent argument below handles every
ordered coefficient triple without importing a homogeneous difference
identity or assuming a reversal symmetry of its forcing sequence.

Other nearby questions have different observables. Vishkautsan studies
strong residual periodicity for products of two reflections on the
unforced Markoff surface, with periodic conics explaining the
local--global phenomenon \citep{vishkautsan2016residual}.
Abboud proves rigidity statements comparing periodic sets of
loxodromic affine-surface automorphisms
\citep[Theorem A]{abboud2025rigidity}.
Planat, Chester and Irwin discuss selected algebraic Painlev\'e VI
solutions in their study of Fricke surfaces
\citep{planat2024dynamics}. These are contextual comparisons, not
dependencies of the present exhaustion proof.

\begin{table}[t]
\centering
\small
\begin{tabularx}{\textwidth}{@{}p{0.20\textwidth}YY@{}}
\toprule
Input or comparison & Scope already available & Scope proved here \\
\midrule
Shin \citep{shin2026character}
& Independent coefficients; height and low-coordinate group-orbit reduction
& Exhaustion for the prescribed three-factor word, uniformly in the level \\
Cantat \citep{cantat2009bers}
& Fixed-fibre dynamics and the resulting integral finiteness
& An explicit residual radius depending only on \(A,B,C\) \\
C421 \citep{c4212026integral}
& The complete equal-forcing scalar classification
& Phase-dependent forcing, including unequal ordered coefficients \\
Elementary algorithms
& Affine returns, finite graphs, polynomial gcds and cycle products
& Their applicability to every remaining point, by the global exhaustion proof \\
\bottomrule
\end{tabularx}
\caption{Scope comparison. The last column is one integrated classification,
not a list of independently new algorithmic ingredients.}
\label{tab:ownership}
\end{table}

The mathematical step is the implication that every periodic orbit
above \(R\) belongs to a finite union of whole periodic lines. We prove
it by a global-maximum entry, forced finite-state propagation and
backward exhaustion using bijective line maps. The source comparisons
above concern the inspected statements and are not a worldwide priority
certificate. Section~\ref{sec:phase} fixes the clock and maximum;
Sections~\ref{sec:graph}--\ref{sec:lines} prove the line exhaustion;
Sections~\ref{sec:effective}--\ref{sec:levels} give the exact residual
and level-count procedures.
